\begin{proposition}[\texorpdfstring{$P_{10}$}{P10} 十次分裂域与 Lee--Yang 三次分裂域的线性互素性]\label{prop:fold-gauge-anomaly-P10-leyang-linear-disjointness}
\leanverified{paper\_fold\_gauge\_anomaly\_P10\_leyang\_linear\_disjointness}
令 \(L/\QQ\) 为命题 \ref{prop:fold-gauge-anomaly-P10-galois-discriminant} 中分歧多项式 \(P_{10}(u)\) 的分裂域，则 \(\mathrm{Gal}(L/\QQ)\cong S_{10}\) 且其唯一指数 \(2\) 子域为
\(
L^{A_{10}}=\QQ(\sqrt{66269989})
\)。
令
\[
P_{\mathrm{LY}}(y):=256y^{3}+411y^{2}+165y+32\in\ZZ[y],
\qquad
L_{\mathrm{LY}}:=\mathrm{Spl}\bigl(P_{\mathrm{LY}}\bigr),
\]
则 \(\mathrm{Gal}(L_{\mathrm{LY}}/\QQ)\cong S_3\)，且其唯一指数 \(2\) 子域为
\(
L_{\mathrm{LY}}^{A_3}=\QQ(\sqrt{-111})
\)
（结论 \ref{con:xi-terminal-zm-leyang-cubic-arithmetic}）。
因此
\[
L\cap L_{\mathrm{LY}}=\QQ,
\qquad\text{从而}\qquad
\mathrm{Gal}(LL_{\mathrm{LY}}/\QQ)\cong S_{10}\times S_{3}.
\]
\end{proposition}
\begin{proof}
由于 \(L/\QQ\) 与 \(L_{\mathrm{LY}}/\QQ\) 均为伽罗瓦扩张，交域 \(L\cap L_{\mathrm{LY}}\) 亦为 \(\QQ\) 上伽罗瓦扩张。
作为 \(L\) 的伽罗瓦子扩张，\(L\cap L_{\mathrm{LY}}\in\{\QQ,\QQ(\sqrt{66269989})\}\)；
作为 \(L_{\mathrm{LY}}\) 的伽罗瓦子扩张，\(L\cap L_{\mathrm{LY}}\in\{\QQ,\QQ(\sqrt{-111})\}\)。
但 \(\QQ(\sqrt{66269989})\) 为实二次域而 \(\QQ(\sqrt{-111})\) 为虚二次域，二者不相等，故只能 \(L\cap L_{\mathrm{LY}}=\QQ\)。
交域平凡且两边皆为伽罗瓦扩张，推出线性互素并得到直积型伽罗瓦群同构。证毕。
\end{proof}

\begin{corollary}[十次分歧域与 Lee--Yang 三次域的联合 Chebotarev 乘法律]\label{cor:fold-gauge-anomaly-P10-leyang-chebotarev-product-law}
\leanverified{paper\_fold\_gauge\_anomaly\_P10\_leyang\_chebotarev\_product\_law}
沿用命题 \ref{prop:fold-gauge-anomaly-P10-leyang-linear-disjointness} 的记号，并写
\[
L_{10}:=L.
\]
对任意满足
\[
p\nmid \mathrm{Disc}(P_{10})\mathrm{Disc}(P_{\mathrm{LY}})
\]
的素数 \(p\)，记
\[
\sigma_{10}(p):=\mathrm{Frob}_p(L_{10}/\QQ)\in S_{10},
\qquad
\sigma_{3}(p):=\mathrm{Frob}_p(L_{\mathrm{LY}}/\QQ)\in S_{3}.
\]
则对任意共轭类
\[
C_{10}\subset S_{10},
\qquad
C_{3}\subset S_{3},
\]
满足
\[
\sigma_{10}(p)\in C_{10},
\qquad
\sigma_{3}(p)\in C_{3}
\]
的素数集合具有自然密度
\[
\boxed{\;
\delta\Bigl(
\sigma_{10}(p)\in C_{10},\ \sigma_{3}(p)\in C_{3}
\Bigr)
=
\frac{|C_{10}|}{|S_{10}|}\cdot\frac{|C_{3}|}{|S_{3}|}.\;}
\]
其中由
\[
L_{10}^{A_{10}}=\QQ(\sqrt{66269989}),
\qquad
L_{\mathrm{LY}}^{A_3}=\QQ(\sqrt{-111})
\]
可见，显式暴露出的二次分歧支撑至少包含四个基本坏素
\[
3,\qquad 37,\qquad 1931,\qquad 34319.
\]
等价地，\(P_{10}\bmod p\) 在 \(\FF_p[u]\) 上的分解型与 \(P_{\mathrm{LY}}\bmod p\) 在 \(\FF_p[y]\) 上的分解型在自然密度意义下严格独立。
\end{corollary}
\begin{proof}[证明要点]
由命题 \ref{prop:fold-gauge-anomaly-P10-leyang-linear-disjointness}，
\[
\mathrm{Gal}(L_{10}L_{\mathrm{LY}}/\QQ)\cong S_{10}\times S_3.
\]
Chebotarev 密度定理说明：对共同不分歧素数，Frobenius 共轭类在该直积群中按共轭类大小等分布。而直积群中的共轭类正是笛卡尔积 \(C_{10}\times C_3\)，其比例为
\[
\frac{|C_{10}|\,|C_3|}{|S_{10}|\,|S_3|}
=
\frac{|C_{10}|}{|S_{10}|}\cdot\frac{|C_3|}{|S_3|}.
\]
对非分歧素数，Frobenius 循环型与模素分解型一致，故得到所示乘法律。证毕。
\end{proof}

\leanverified{paper\_fold\_gauge\_anomaly\_p10\_leyang\_irreducible\_root\_split\_densities}
\begin{corollary}[十次不可约与 Lee--Yang 单根/全分裂的显式正密度]\label{cor:fold-gauge-anomaly-P10-leyang-irreducible-root-split-densities}
\leanverified{paper\_fold\_gauge\_anomaly\_p10\_leyang\_irreducible\_root\_split\_densities}
在推论 \ref{cor:fold-gauge-anomaly-P10-leyang-chebotarev-product-law} 的设定下，记
\[
g(y):=P_{\mathrm{LY}}(y)=256y^{3}+411y^{2}+165y+32.
\]
对共同不分歧素数 \(p\)，定义事件
\[
\mathcal I:=\{\text{\(P_{10}\bmod p\) 在 \(\FF_p[u]\) 上不可约}\},
\]
\[
\mathcal R:=\{\text{\(g\bmod p\) 在 \(\FF_p[y]\) 上恰有一个 \(\FF_p\)-根}\},
\qquad
\mathcal S:=\{\text{\(g\bmod p\) 在 \(\FF_p[y]\) 上完全分裂}\}.
\]
则有显式正密度
\[
\boxed{\;
\delta(\mathcal I\cap \mathcal R)
=
\frac{1}{10}\cdot\frac12
=
\frac1{20}.\;}
\]
\[
\boxed{\;
\delta(\mathcal I\cap \mathcal S)
=
\frac{1}{10}\cdot\frac16
=
\frac1{60}.\;}
\]
\end{corollary}
\begin{proof}
由推论 \ref{cor:fold-gauge-anomaly-P10-modp-splitting-densities}，事件 \(\mathcal I\) 对应 \(S_{10}\) 中的 \(10\)-循环共轭类，故
\[
\delta(\mathcal I)=\frac1{10}.
\]
另一方面，由结论 \ref{con:xi-terminal-zm-leyang-cubic-modp-splitting-density}，\(g\bmod p\) 的分解型与 \(S_3\) 的循环型一一对应：换位类对应 \((1)(2)\)，故
\[
\delta(\mathcal R)=\frac{3}{6}=\frac12;
\]
恒等类对应 \((1)(1)(1)\)，故
\[
\delta(\mathcal S)=\frac{1}{6}.
\]
将这些边际密度代入推论 \ref{cor:fold-gauge-anomaly-P10-leyang-chebotarev-product-law} 的乘法律，即得两式。证毕。
\end{proof}

\begin{corollary}[最大阿贝尔子扩张：双二次判别角色]\label{cor:fold-gauge-anomaly-P10-leyang-max-abelian-subextension}
\leanverified{paper\_fold\_gauge\_anomaly\_p10\_leyang\_max\_abelian\_subextension}
令 \(F:=LL_{\mathrm{LY}}\)。
则 \(F/\QQ\) 的最大阿贝尔子扩张为
\[
F^{\mathrm{ab}}
=\QQ\!\left(\sqrt{66269989},\sqrt{-111}\right),
\qquad
\mathrm{Gal}(F^{\mathrm{ab}}/\QQ)\cong (\ZZ/2\ZZ)^2.
\]
并且对 \(p\nmid 3\cdot 37\cdot 1931\cdot 34319\) 的素数，\(p\) 在 \(F^{\mathrm{ab}}\) 中的分解行为由两条二次特征
\[
\left(\frac{66269989}{p}\right),\qquad \left(\frac{-111}{p}\right)
\]
完全刻画；四种符号组合在自然密度意义下各占 \(1/4\)。
\end{corollary}
\begin{proof}
由命题 \ref{prop:fold-gauge-anomaly-P10-leyang-linear-disjointness} 得 \(\mathrm{Gal}(F/\QQ)\cong S_{10}\times S_3\)。
而 \(S_{10}^{\mathrm{ab}}\cong C_2\)、\(S_3^{\mathrm{ab}}\cong C_2\)，故
\(
(S_{10}\times S_3)^{\mathrm{ab}}\cong C_2\times C_2
\)；
对应的最大阿贝尔子扩张由两因子的唯一二次子域生成，即为所示双二次域。
最后的 \(1/4\) 等分布由双二次扩张的 Chebotarev 定理直接推出。证毕。
\end{proof}

\begin{theorem}[Frobenius 类函数统计的严格张量独立]\label{thm:fold-gauge-anomaly-P10-leyang-class-function-tensor-independence}
沿用命题 \ref{prop:fold-gauge-anomaly-P10-leyang-linear-disjointness} 的记号。对任意类函数
\[
f:S_{10}\to\CC,
\qquad
g:S_3\to\CC,
\]
记群平均
\[
\langle f\rangle_{S_{10}}:=\frac{1}{10!}\sum_{\tau\in S_{10}}f(\tau),
\qquad
\langle g\rangle_{S_3}:=\frac{1}{6}\sum_{\upsilon\in S_3}g(\upsilon).
\]
再记
\[
\pi_{10,\mathrm{LY}}^\ast(x):=
\#\left\{p\le x:\ p\nmid \Disc(P_{10})\Disc(P_{\mathrm{LY}})\right\}.
\]
则有极限恒等式
\[
\boxed{\;
\lim_{x\to\infty}
\frac{1}{\pi_{10,\mathrm{LY}}^\ast(x)}
\sum_{\substack{p\le x\\ p\nmid \Disc(P_{10})\Disc(P_{\mathrm{LY}})}}
f\!\bigl(\sigma_{10}(p)\bigr)\,
g\!\bigl(\sigma_{3}(p)\bigr)
=
\langle f\rangle_{S_{10}}\,
\langle g\rangle_{S_3}.\;}
\]
特别地，若 \(\chi^\lambda\) 与 \(\psi^\mu\) 分别为 \(S_{10}\) 与 \(S_3\) 的不可约特征标，则
\[
\boxed{\;
\lim_{x\to\infty}
\frac{1}{\pi_{10,\mathrm{LY}}^\ast(x)}
\sum_{\substack{p\le x\\ p\nmid \Disc(P_{10})\Disc(P_{\mathrm{LY}})}}
\chi^\lambda\!\bigl(\sigma_{10}(p)\bigr)\,
\psi^\mu\!\bigl(\sigma_3(p)\bigr)
=
\begin{cases}
1,& \chi^\lambda=\mathbf 1,\ \psi^\mu=\mathbf 1,\\
0,& \text{否则}.
\end{cases}\;}
\]
\end{theorem}
\leanverified{paper\_fold\_gauge\_anomaly\_p10\_leyang\_class\_function\_tensor\_independence}
\begin{proof}
把 \(f\) 与 \(g\) 分别写成共轭类指示函数的线性组合：
\[
f=\sum_C a_C\,\mathbf 1_C,
\qquad
g=\sum_D b_D\,\mathbf 1_D,
\]
其中 \(C\subset S_{10}\)、\(D\subset S_3\) 遍历共轭类。于是
\[
f\!\bigl(\sigma_{10}(p)\bigr)\,g\!\bigl(\sigma_3(p)\bigr)
=
\sum_{C,D}a_Cb_D\,
\mathbf 1_{\{\sigma_{10}(p)\in C,\ \sigma_3(p)\in D\}}.
\]
对非分歧素数取平均，并应用推论 \ref{cor:fold-gauge-anomaly-P10-leyang-chebotarev-product-law}，得到
\[
\lim_{x\to\infty}
\frac{1}{\pi_{10,\mathrm{LY}}^\ast(x)}
\sum_{\substack{p\le x\\ p\nmid \Disc(P_{10})\Disc(P_{\mathrm{LY}})}}
f\!\bigl(\sigma_{10}(p)\bigr)\,g\!\bigl(\sigma_3(p)\bigr)
=
\sum_{C,D}a_Cb_D\frac{|C|}{10!}\frac{|D|}{6}.
\]
右端按 \(C\) 与 \(D\) 分组即化为
\[
\left(\sum_C a_C\frac{|C|}{10!}\right)
\left(\sum_D b_D\frac{|D|}{6}\right)
=
\langle f\rangle_{S_{10}}\langle g\rangle_{S_3}.
\]
这就证明了第一式。

对不可约特征标，使用有限群的标准平均恒等式
\[
\frac{1}{|G|}\sum_{g\in G}\chi(g)=
\begin{cases}
1,& \chi=\mathbf 1,\\
0,& \chi\neq \mathbf 1,
\end{cases}
\]
分别作用于 \(G=S_{10}\) 与 \(G=S_3\)，即可得到第二式。证毕。
\end{proof}

\begin{corollary}[十次判别式符号与 Lee--Yang 三次分解型的零互信息]\label{cor:fold-gauge-anomaly-P10-leyang-zero-mutual-information}
\leanverified{paper\_fold\_gauge\_anomaly\_p10\_leyang\_zero\_mutual\_information}
对不分歧素数 \(p\)，定义
\[
\chi_{10}(p):=\mathrm{sgn}\!\bigl(\sigma_{10}(p)\bigr)
=
\left(\frac{66269989}{p}\right)\in\{\pm 1\},
\]
并记 \(\mathcal T_{\mathrm{LY}}(p)\) 为 \(P_{\mathrm{LY}}\bmod p\) 的分解型随机变量，其取值于
\[
\{(1)(1)(1),\ (1)(2),\ (3)\}.
\]
将不分歧素数集合赋予自然密度概率测度，则 \(\chi_{10}\) 与 \(\mathcal T_{\mathrm{LY}}\) 严格独立。等价地，对任意
\[
\tau\in\{(1)(1)(1),(1)(2),(3)\},
\qquad
\varepsilon\in\{\pm 1\},
\]
都有
\[
\boxed{\;
\mathbb{P}\!\bigl(\mathcal T_{\mathrm{LY}}=\tau,\ \chi_{10}=\varepsilon\bigr)
=
\mathbb{P}\!\bigl(\mathcal T_{\mathrm{LY}}=\tau\bigr)\,
\mathbb{P}\!\bigl(\chi_{10}=\varepsilon\bigr).\;}
\]
从而
\[
\boxed{\;
I\!\bigl(\chi_{10};\mathcal T_{\mathrm{LY}}\bigr)=0.\;}
\]
更一般地，对任意函数
\[
h:\{(1)(1)(1),(1)(2),(3)\}\to\CC
\]
都有
\[
\boxed{\;
\mathbb{E}\!\left[h(\mathcal T_{\mathrm{LY}})\mid \chi_{10}=1\right]
=
\mathbb{E}\!\left[h(\mathcal T_{\mathrm{LY}})\mid \chi_{10}=-1\right]
=
\frac16 h((1)(1)(1))+\frac12 h((1)(2))+\frac13 h((3)).\;}
\]
\end{corollary}
\begin{proof}
取 \(f_\varepsilon:S_{10}\to\{0,1\}\) 为符号类函数
\[
f_\varepsilon(\tau):=\mathbf 1_{\{\mathrm{sgn}(\tau)=\varepsilon\}},
\]
并取 \(g_\tau:S_3\to\{0,1\}\) 为分解型 \(\tau\) 所对应的共轭类指示函数。由定理
\ref{thm:fold-gauge-anomaly-P10-leyang-class-function-tensor-independence}，
\[
\mathbb{P}\!\bigl(\mathcal T_{\mathrm{LY}}=\tau,\ \chi_{10}=\varepsilon\bigr)
=
\langle f_\varepsilon\rangle_{S_{10}}\,
\langle g_\tau\rangle_{S_3}
=
\mathbb{P}\!\bigl(\chi_{10}=\varepsilon\bigr)\,
\mathbb{P}\!\bigl(\mathcal T_{\mathrm{LY}}=\tau\bigr).
\]
故两随机变量独立。

独立的离散随机变量互信息为零，得第二式。最后，条件期望公式只是把独立性代入
\[
\mathbb{E}\!\left[h(\mathcal T_{\mathrm{LY}})\mid \chi_{10}=\varepsilon\right]
=
\sum_{\tau}h(\tau)\,
\mathbb{P}\!\bigl(\mathcal T_{\mathrm{LY}}=\tau\mid \chi_{10}=\varepsilon\bigr)
\]
并使用 \(S_3\) 三个共轭类大小为 \(1,3,2\)，故其边际分布分别为 \(1/6,1/2,1/3\)。证毕。
\end{proof}

\begin{proposition}[\texorpdfstring{$P_{10}$}{P10} 十次分裂域与指数平方因子 \texorpdfstring{$H$}{H} 的线性互素性]\label{prop:fold-gauge-anomaly-P10-H-linear-disjointness}
令 \(L/\QQ\) 为命题 \ref{prop:fold-gauge-anomaly-P10-galois-discriminant} 中 \(P_{10}(u)\) 的分裂域。
令 \(H(u)\) 为命题 \ref{prop:fold-gauge-anomaly-H-galois-discriminant} 中次数 \(19\) 的不可约多项式，并记其分裂域为 \(L_H:=\mathrm{Spl}(H)\)。
则
\[
L\cap L_H=\QQ,
\qquad\text{从而}\qquad
\mathrm{Gal}(LL_H/\QQ)\cong S_{10}\times S_{19}.
\]
\end{proposition}
\begin{proof}
记 \(E:=L\cap L_H\)。
由于 \(L/\QQ\) 与 \(L_H/\QQ\) 均为伽罗瓦扩张，\(E/\QQ\) 亦为伽罗瓦扩张。
\par
由命题 \ref{prop:fold-gauge-anomaly-P10-galois-discriminant}，\(L\) 的唯一指数 \(2\) 子域为
\(
K_{10}:=L^{A_{10}}=\QQ(\sqrt{66269989})
\)，故 \(E\in\{\QQ,K_{10}\}\)。
由命题 \ref{prop:fold-gauge-anomaly-H-galois-discriminant}，\(\mathrm{Gal}(L_H/\QQ)\cong S_{19}\) 且其唯一指数 \(2\) 子域为
\[
K_H:=L_H^{A_{19}}=\QQ\!\left(\sqrt{\mathrm{Disc}(H)}\right),
\]
故 \(E\in\{\QQ,K_H\}\)。
\par
由命题 \ref{prop:fold-gauge-anomaly-H-galois-discriminant} 的判别式分解可见素数 \(101\) 在 \(\mathrm{Disc}(H)\) 中以奇指数出现，因而 \(101\) 在二次域 \(K_H\) 中分歧。
而 \(101\nmid 66269989\)，故 \(101\) 在 \(K_{10}\) 中不分歧，从而 \(K_{10}\neq K_H\)。
于是 \(E\) 不可能为任何非平凡二次子域，必有 \(E=\QQ\)。
交域平凡且两边皆为伽罗瓦扩张，推出线性互素并得到直积型伽罗瓦群同构。证毕。
\leanverified{paper\_fold\_gauge\_anomaly\_P10\_H\_linear\_disjointness}
\leanverified{paper\_fold\_gauge\_anomaly\_p10\_h\_linear\_disjointness}
\end{proof}

\begin{corollary}[合成分裂域的直积分解]\label{cor:fold-gauge-anomaly-P10H-splitting-field-product}
记 \(D(u):=P_{10}(u)\,H(u)\in\ZZ[u]\)。
则 \(D\) 的分裂域满足
\[
\mathrm{Spl}(D)=LL_H,
\qquad
\mathrm{Gal}\bigl(\mathrm{Spl}(D)/\QQ\bigr)\cong S_{10}\times S_{19}.
\]
\end{corollary}
\begin{proof}
\(\mathrm{Spl}(D)\) 由 \(P_{10}\) 与 \(H\) 的根集合共同生成，故 \(\mathrm{Spl}(D)=LL_H\)。
伽罗瓦群同构由命题 \ref{prop:fold-gauge-anomaly-P10-H-linear-disjointness} 直接推出。证毕。
\end{proof}
\leanverified{paper\_fold\_gauge\_anomaly\_P10H\_splitting\_field\_product}
\leanverified{paper\_fold\_gauge\_anomaly\_p10h\_splitting\_field\_product}

\begin{corollary}[Chebotarev 乘积律：\texorpdfstring{$P_{10}$}{P10} 与 \texorpdfstring{$H$}{H} 的分解型完全解耦]\label{cor:fold-gauge-anomaly-P10-H-chebotarev-product-law}
\leanverified{paper\_fold\_gauge\_anomaly\_P10\_H\_chebotarev\_product\_law}
在命题 \ref{prop:fold-gauge-anomaly-P10-H-linear-disjointness} 的记号下，取素数 \(p\nmid \mathrm{Disc}(P_{10})\mathrm{Disc}(H)\)。
令 \(\lambda\vdash 10\)、\(\mu\vdash 19\) 为分拆，并记 \(C_\lambda\subset S_{10}\)、\(C_\mu\subset S_{19}\) 为相应循环型的共轭类。
则满足 \(P_{10}\bmod p\) 的分解型为 \(\lambda\) 且 \(H\bmod p\) 的分解型为 \(\mu\) 的素数集合具有自然密度
\[
\delta(\lambda,\mu)=\frac{|C_\lambda|}{10!}\cdot \frac{|C_\mu|}{19!}.
\]
\end{corollary}
\begin{proof}
由命题 \ref{prop:fold-gauge-anomaly-P10-H-linear-disjointness}，\(\mathrm{Gal}(LL_H/\QQ)\cong S_{10}\times S_{19}\)。
对非分歧素数 \(p\)，Frobenius 共轭类落在 \(C_\lambda\times C_\mu\) 中当且仅当两边模素分解型分别为 \(\lambda\) 与 \(\mu\)。
Chebotarev 密度定理给出密度为
\(
|C_\lambda\times C_\mu|/|S_{10}\times S_{19}|
=\frac{|C_\lambda|}{10!}\cdot\frac{|C_\mu|}{19!}
\)。
证毕。
\end{proof}

\begin{corollary}[若干可核验的显式密度常数]\label{cor:fold-gauge-anomaly-P10-H-explicit-densities}
\leanverified{paper\_fold\_gauge\_anomaly\_p10\_h\_explicit\_densities}
在推论 \ref{cor:fold-gauge-anomaly-P10-H-chebotarev-product-law} 的条件下，以下密度均指除去有限个分歧素数后的自然密度：
\begin{enumerate}
  \item 同时不可约：
  \[
  \delta\bigl((10),(19)\bigr)=\frac1{10}\cdot\frac1{19}=\frac1{190}.
  \]
  \item \(P_{10}\) 含一次因子且 \(H\) 不可约：
  \[
  \delta\bigl(\text{$P_{10}$ 含一次因子},(19)\bigr)
  =
  \Bigl(1-\frac{!10}{10!}\Bigr)\cdot\frac1{19}
  =
  \frac{28319}{44800}\cdot\frac1{19}.
  \]
  \item \(P_{10}\) 与 \(H\) 同时含一次因子：
  \[
  \delta\bigl(\text{$P_{10}$ 含一次因子},\text{$H$ 含一次因子}\bigr)
  =
  \Bigl(1-\frac{!10}{10!}\Bigr)\cdot \Bigl(1-\frac{!19}{19!}\Bigr).
  \]
\end{enumerate}
\end{corollary}
\begin{proof}
由推论 \ref{cor:fold-gauge-anomaly-P10-H-chebotarev-product-law}，密度分解为边际密度之乘积。
\(\delta((10))=1/10\)、\(\delta((19))=1/19\)。
并且 \(\delta(\text{$P_{10}$ 含一次因子})=1-!10/10!=28319/44800\) 由推论 \ref{cor:fold-gauge-anomaly-P10-modp-splitting-densities}。
而 \(\delta(\text{$H$ 含一次因子})=1-!19/19!\) 为 \(S_{19}\) 中存在不动点的比例。
证毕。
\end{proof}

\begin{corollary}[最大阿贝尔子扩张与双二次特征的等分布]\label{cor:fold-gauge-anomaly-P10-H-max-abelian-and-quadratic-equidistribution}
\leanverified{paper\_fold\_gauge\_anomaly\_p10\_h\_max\_abelian\_and\_quadratic\_equidistribution}
设 \(F:=LL_H\)。
则 \(F/\QQ\) 的最大阿贝尔子扩张为
\[
F^{\mathrm{ab}}=K_{10}K_H,
\qquad
\mathrm{Gal}(F^{\mathrm{ab}}/\QQ)\cong (\ZZ/2\ZZ)^2,
\]
其中 \(K_{10}=\QQ(\sqrt{66269989})\)、\(K_H=\QQ(\sqrt{\mathrm{Disc}(H)})\)。
并且对 \(p\) 不在 \(K_{10}K_H\) 中分歧的素数，二次特征对
\[
\chi_{10}(p):=\left(\frac{66269989}{p}\right),
\qquad
\chi_H(p):=\left(\frac{\mathrm{Disc}(H)}{p}\right)
\]
的四种符号组合在自然密度意义下各占 \(1/4\)。
\end{corollary}
\begin{proof}
由命题 \ref{prop:fold-gauge-anomaly-P10-H-linear-disjointness}，
\(\mathrm{Gal}(F/\QQ)\cong S_{10}\times S_{19}\)。
而 \(S_{10}^{\mathrm{ab}}\cong C_2\)、\(S_{19}^{\mathrm{ab}}\cong C_2\)，故
\(
(S_{10}\times S_{19})^{\mathrm{ab}}\cong C_2\times C_2
\)；
对应的最大阿贝尔子扩张由两因子的唯一二次子域生成，即 \(K_{10}K_H\)。
对双二次扩张应用 Chebotarev 定理得到 Frobenius 在 \((\ZZ/2\ZZ)^2\) 上等分布，从而四种符号组合密度均为 \(1/4\)。证毕。
\end{proof}

\begin{proposition}[固定一次因子数的 Rencontres 密度律]\label{prop:fold-gauge-anomaly-H-rencontres-fixed-point-density}
\leanverified{paper\_fold\_gauge\_anomaly\_H\_rencontres\_fixed\_point\_density}
\leanverified{paper\_fold\_gauge\_anomaly\_h\_rencontres\_fixed\_point\_density}
对 \(r\in\{0,1,\dots,19\}\)，记事件
\[
\mathcal E_r(p):=\text{$H\bmod p$ 恰有 \(r\) 个一次因子}.
\]
则除去有限个分歧素数后，\(\mathcal E_r\) 的自然密度为
\[
\delta(\mathcal E_r)
=\frac{\binom{19}{r}\,!(19-r)}{19!}
=\frac{1}{r!}\sum_{k=0}^{19-r}\frac{(-1)^k}{k!},
\]
其中 \(!(m)\) 为 \(m\) 阶错排数。
进一步，若记事件 \(\mathcal I(p):=\text{$P_{10}\bmod p$ 不可约}\)，则有联合密度
\[
\delta(\mathcal I\cap \mathcal E_r)=\frac1{10}\cdot \delta(\mathcal E_r).
\]
\end{proposition}
\begin{proof}
由命题 \ref{prop:fold-gauge-anomaly-H-galois-discriminant}，\(\mathrm{Gal}(H/\QQ)\cong S_{19}\)。
对非分歧素数 \(p\)，\(H\bmod p\) 的一次因子个数等于 Frobenius 置换在 \(19\) 个根上不动点数。
恰有 \(r\) 个不动点的置换个数为 \(\binom{19}{r}!(19-r)\)，故密度为 \(\binom{19}{r}!(19-r)/19!\)。
等价表达式由错排数的容斥公式化简得到。
\par
联合密度由推论 \ref{cor:fold-gauge-anomaly-P10-H-chebotarev-product-law} 的直积独立性给出：\(\delta(\mathcal I)=1/10\) 且与 \(\mathcal E_r\) 独立，故
\(
\delta(\mathcal I\cap\mathcal E_r)=\delta(\mathcal I)\delta(\mathcal E_r)=\frac1{10}\delta(\mathcal E_r)
\)。
证毕。
\end{proof}

\begin{theorem}[Lee--Yang 三次域、\texorpdfstring{$P_{10}$}{P10} 十次域与指数域的三重直积解耦]\label{thm:fold-gauge-anomaly-P10-leyang-H-triple-product}
\leanverified{paper\_fold\_gauge\_anomaly\_p10\_leyang\_h\_triple\_product}
在命题 \ref{prop:fold-gauge-anomaly-P10-leyang-linear-disjointness} 与命题 \ref{prop:fold-gauge-anomaly-P10-H-linear-disjointness} 的记号下，有
\[
\mathrm{Gal}(LL_{\mathrm{LY}}L_H/\QQ)\cong S_{10}\times S_{3}\times S_{19}.
\]
\end{theorem}
\begin{proof}
记 \(F:=LL_{\mathrm{LY}}\) 并由命题 \ref{prop:fold-gauge-anomaly-P10-leyang-linear-disjointness} 得 \(\mathrm{Gal}(F/\QQ)\cong S_{10}\times S_3\)。
由推论 \ref{cor:fold-gauge-anomaly-P10-leyang-max-abelian-subextension}，\(F\) 的最大阿贝尔子扩张为
\(
F^{\mathrm{ab}}=\QQ(\sqrt{66269989},\sqrt{-111})
\)。
\par
设 \(E:=F\cap L_H\)。
由于 \(\mathrm{Gal}(L_H/\QQ)\cong S_{19}\)，其伽罗瓦子扩张只有 \(\QQ\) 与唯一二次子域 \(K_H=L_H^{A_{19}}\)，故 \(E\in\{\QQ,K_H\}\)。
由命题 \ref{prop:fold-gauge-anomaly-H-galois-discriminant}，\(101\) 在 \(K_H\) 中分歧。
但 \(101\) 在 \(F^{\mathrm{ab}}\) 中不分歧（其分歧素数只可能来自 \(66269989\) 与 \(-111\) 的素因子），故 \(K_H\not\subseteq F\)，从而 \(E=\QQ\)。
于是 \(F\) 与 \(L_H\) 线性互素并得到
\(
\mathrm{Gal}(FL_H/\QQ)\cong \mathrm{Gal}(F/\QQ)\times \mathrm{Gal}(L_H/\QQ)
\cong S_{10}\times S_3\times S_{19}
\)。
证毕。
\end{proof}
\leanverified{paper\_fold\_gauge\_anomaly\_P10\_leyang\_H\_triple\_product}

\begin{corollary}[三重分解型的完全乘法密度]\label{cor:fold-gauge-anomaly-P10-leyang-H-chebotarev-triple-product}
在定理 \ref{thm:fold-gauge-anomaly-P10-leyang-H-triple-product} 的设定下，
对任意共轭类 \(C_\lambda\subset S_{10}\)、\(C_\nu\subset S_3\)、\(C_\mu\subset S_{19}\)，满足 \(\mathrm{Frob}_p\in C_\lambda\times C_\nu\times C_\mu\) 的素数集合（除去有限坏素集合后）具有自然密度
\[
\delta(\lambda,\nu,\mu)=\frac{|C_\lambda|}{10!}\cdot\frac{|C_\nu|}{3!}\cdot\frac{|C_\mu|}{19!}.
\]
\end{corollary}
\begin{proof}
由定理 \ref{thm:fold-gauge-anomaly-P10-leyang-H-triple-product}，伽罗瓦群为直积。
对直积群应用 Chebotarev 密度定理即可得到共轭类分布的乘积分解。证毕。
\end{proof}
\leanverified{paper\_fold\_gauge\_anomaly\_P10\_leyang\_H\_chebotarev\_triple\_product}
\leanverified{paper\_fold\_gauge\_anomaly\_p10\_leyang\_h\_chebotarev\_triple\_product}

\begin{proposition}[Lee--Yang 三次分裂域与 \texorpdfstring{$\ell$}{l}-挠域的线性互素性]\label{prop:fold-gauge-anomaly-leyang-elliptic-torsion-disjointness}
\leanverified{paper\_fold\_gauge\_anomaly\_leyang\_elliptic\_torsion\_disjointness}
沿用命题 \ref{prop:fold-gauge-anomaly-P10-elliptic-torsion-disjointness} 的椭圆曲线 \(E/\QQ\) 与挠域 \(M_\ell=\QQ(E[\ell])\) 记号，并记 \(L_{\mathrm{LY}}=\mathrm{Spl}(P_{\mathrm{LY}})\)。
则对任意素数 \(\ell\neq 3\)，有
\[
L_{\mathrm{LY}}\cap M_\ell=\QQ,
\qquad\text{从而}\qquad
\mathrm{Gal}(L_{\mathrm{LY}}M_\ell/\QQ)\cong S_{3}\times \mathrm{Gal}(M_\ell/\QQ).
\]
\end{proposition}
\begin{proof}
交域 \(L_{\mathrm{LY}}\cap M_\ell\) 为 \(\QQ\) 上伽罗瓦扩张，且作为 \(L_{\mathrm{LY}}\) 的伽罗瓦子扩张只能为 \(\QQ\) 或其唯一二次子域 \(\QQ(\sqrt{-111})\)（结论 \ref{con:xi-terminal-zm-leyang-cubic-arithmetic}）。
但 \(\QQ(\sqrt{-111})\) 在素数 \(3\) 处分歧，而当 \(\ell\neq 3\) 时，N\'eron--Ogg--Shafarevich 判据给出 \(M_\ell/\QQ\) 仅可能在 \(37\) 与 \(\ell\) 处分歧（命题 \ref{prop:fold-gauge-anomaly-P10-elliptic-torsion-disjointness}），故其任意子扩张在 \(3\) 处亦不分歧。
矛盾，遂交域只能为 \(\QQ\)；两边皆为伽罗瓦扩张且交域平凡，得到线性互素与直积型伽罗瓦群同构。证毕。
\end{proof}

\begin{corollary}[三通道直积独立性：\texorpdfstring{$P_{10}$}{P10}、Lee--Yang 与 \texorpdfstring{$\ell$}{l}-挠域]\label{cor:fold-gauge-anomaly-P10-leyang-elliptic-triple-independence}
\leanverified{paper\_fold\_gauge\_anomaly\_p10\_leyang\_elliptic\_triple\_independence}
在命题 \ref{prop:fold-gauge-anomaly-P10-galois-discriminant}、命题 \ref{prop:fold-gauge-anomaly-P10-leyang-linear-disjointness}、命题 \ref{prop:fold-gauge-anomaly-P10-elliptic-torsion-disjointness} 与命题 \ref{prop:fold-gauge-anomaly-leyang-elliptic-torsion-disjointness} 的记号下，
若 \(\ell\notin\{3,1931,34319\}\)，则
\[
\mathrm{Gal}(LL_{\mathrm{LY}}M_\ell/\QQ)
\ \cong\
S_{10}\times S_{3}\times \mathrm{Gal}(M_\ell/\QQ).
\]
因此对除去有限坏素集合后的素数 \(p\)，下列三类数据在自然密度意义下相互独立：
\begin{enumerate}
  \item \(P_{10}\bmod p\) 的分解型（对应 \(S_{10}\) 的 Frobenius 循环型）；
  \item \(P_{\mathrm{LY}}\bmod p\) 的分解型（对应 \(S_{3}\) 的 Frobenius 循环型）；
  \item \(\rho_{E,\ell}(\mathrm{Frob}_p)\) 在 \(\mathrm{Gal}(M_\ell/\QQ)\) 中的共轭类。
\end{enumerate}
\end{corollary}
\begin{proof}[证明要点]
令 \(F:=LL_{\mathrm{LY}}\) 并由命题 \ref{prop:fold-gauge-anomaly-P10-leyang-linear-disjointness} 得 \(\mathrm{Gal}(F/\QQ)\cong S_{10}\times S_3\)。
设 \(K_1=\QQ(\sqrt{66269989})\)、\(K_2=\QQ(\sqrt{-111})\)。
由推论 \ref{cor:fold-gauge-anomaly-P10-leyang-max-abelian-subextension}，\(F\) 的全部二次子域恰为 \(K_1\)、\(K_2\) 与 \(\QQ(\sqrt{66269989\cdot(-111)})\) 三者。
\par
当 \(\ell\notin\{3,1931,34319\}\) 时，N\'eron--Ogg--Shafarevich 判据表明 \(M_\ell/\QQ\) 在 \(3,1931,34319\) 处不分歧（命题 \ref{prop:fold-gauge-anomaly-P10-elliptic-torsion-disjointness}），从而 \(M_\ell\) 不可能包含上述任一二次子域。
另一方面，由于 \(S_{10}\) 与 \(S_3\) 的全部非平凡正规商都经由符号特征产生的 \(C_2\)（命题 \ref{prop:fold-gauge-anomaly-P10-galois-discriminant} 与结论 \ref{con:xi-terminal-zm-leyang-cubic-arithmetic}），任意非平凡伽罗瓦子扩张 \(E/\QQ\) 满足 \(E\subseteq F\) 时必含某一二次子域。
据此 \(F\cap M_\ell\) 只能为 \(\QQ\)，从而 \(F\) 与 \(M_\ell\) 线性互素，给出所示直积同构。
独立性为 Chebotarev 密度定理在直积群上的乘积分布之直接推论。证毕。
\end{proof}

\paragraph{有限域审计接口}
脚本 \texttt{scripts/exp\_fold\_gauge\_anomaly\_s4\_hurwitz\_tower\_audit.py} 在同一 JSON 证书中并列给出 \(P_{10}\) 与 \(P_{\mathrm{LY}}\) 的判别式平方自由核、二次子域分离，以及对一段有限素数范围的联合分解型计数（输出：\texttt{artifacts/export/fold\_gauge\_anomaly\_s4\_hurwitz\_tower\_audit.json}）。

\endinput
